% Supplement prepared by the Part G agent for the parent's B.4 proof.
% No claim number is assigned here; integrate into the numbered B statement.
\paragraph{The alternating simplex is a ball.}
Let $K$ be the nerve of the indiscrete groupoid on the two objects
$+,-$, and let $K_n$ be its $n$-skeleton.  Its nondegenerate
$j$-simplices are the two alternating strings of $j+1$ signs.
We prove that $|K_n|$ is an $n$-sphere, and that the image in $K$
of either alternating $n$-simplex realizes as a closed $n$-ball
whose boundary is $|K_{n-1}|$.  In particular any simplicial subset
of that image which omits its top nondegenerate simplex lies in
the boundary of this ball.

Represent a simplex with sign string $(\epsilon_0,\ldots,\epsilon_j)$
and barycentric coordinates $(\lambda_0,\ldots,\lambda_j)$ by the
step function on $[0,1]$ which has value $\epsilon_i\in\{-1,1\}$
on an interval of length $\lambda_i$, the intervals being placed in
the indicated order.  Deleting a zero coordinate and merging
consecutive equal signs are precisely the realization identifications.
It follows that $|K_n|$ identifies bijectively with the almost
everywhere equivalence classes of $\{-1,1\}$-valued step functions
having at most $n$ sign changes in $(0,1)$.  The map to this space
with its $L^1$ metric is continuous on each simplex: moving the
breakpoints changes the function only on intervals whose total
length tends to zero.  Since $K_n$ has finitely many nondegenerate
simplices, its realization is compact.  The $L^1$ target is Hausdorff,
so this bijection is a homeomorphism.

We give a finite-dimensional convex proof of the asserted sphere and
ball descriptions.  For $d\geq1$, put
\[
 p_c(t)=\sum_{j=0}^{d-1}c_jt^j,
 \qquad h_d(c)=\int_0^1|p_c(t)|\,dt,
 \qquad
 M_d=\{x\in\mathbb R^d:c\cdot x\leq h_d(c)
                    \text{ for every }c\in\mathbb R^d\}.
\]
This is a compact convex body with $0$ in its interior.  Indeed the
coordinate inequalities bound it, and the positive continuous function
$h_d$ has a positive minimum on the Euclidean unit sphere, so a small
ball about $0$ lies in $M_d$.  For nonzero $c$, define
\[
 x(c)_j=\int_0^1t^j\operatorname{sign}(p_c(t))\,dt
       \quad(0\leq j<d).
\]
The polynomial has finitely many zeros.  For every $a$ we have
$a\cdot x(c)\leq h_d(a)$, while
$c\cdot x(c)=h_d(c)$.  Thus $h_d$ is exactly the support function
of $M_d$, and $x(c)$ belongs to its support face with normal $c$.

That support face has only one point.  Dominated convergence gives,
for every vector $a$,
\[
 \left.\frac{d}{ds}\right|_{s=0}h_d(c+sa)
   =\int_0^1p_a(t)\operatorname{sign}(p_c(t))\,dt
   =a\cdot x(c).
\]
If $y$ belongs to the same support face, the inequalities
$h_d(c+sa)\geq(c+sa)\cdot y$ for positive and negative $s$
force $a\cdot y=a\cdot x(c)$ for every $a$, hence $y=x(c)$.
Every boundary point has a supporting hyperplane, so every boundary
point of $M_d$ is the moment vector of the sign of a polynomial of
degree at most $d-1$.

Take $d=n+1$.  Any step function $f$ with at most $n$ sign changes
is, almost everywhere, the sign of a nonzero polynomial of degree
at most $n$: use the product of the linear factors at its actual
change points and choose the sign of the leading coefficient.
Consequently its moment vector
\[
 \mu_n(f)=\left(\int_0^1t^j f(t)\,dt\right)_{0\leq j\leq n}
\]
lies on $\partial M_{n+1}$, and the preceding support argument proves
surjectivity onto that boundary.  This map is also injective.
If $f=\operatorname{sign}(p)$ and $\mu_n(g)=\mu_n(f)$, then
\[
 \int_0^1p(t)g(t)\,dt=\int_0^1|p(t)|\,dt.
\]
The integrand $|p|-pg$ is nonnegative, so equality implies $g=f$
almost everywhere.  Continuity and compactness therefore identify
$|K_n|$ homeomorphically with $\partial M_{n+1}$, which is an
$n$-sphere by radial projection from $0$.

It remains to identify the closure of either top simplex, rather
than merely the whole skeleton.  For $n\geq1$, project $M_{n+1}$
onto its first $n$ coordinates.  Its image is $M_n$: their support
functions agree, because $h_{n+1}(c,0)=h_n(c)$.  Write its fibers
as intervals
\[
 \{x\}\times[\ell(x),u(x)],\qquad x\in M_n.
\]
For $x\in\operatorname{int}(M_n)$ the interval has positive length.
For example, since $0\in\operatorname{int}(M_{n+1})$, express $x$
as a strict convex combination of $0$ and another point of $M_n$;
combining a lift of the latter point with a small ball about $0$
gives a nontrivial vertical interval over $x$.  A support normal
to $(x,u(x))$ must then have positive last coefficient: a negative
coefficient is contradicted by the lower point in the same fiber,
and a zero coefficient would support $M_n$ at its interior point
$x$.  Its polynomial therefore has degree $n$ and positive leading
coefficient.  Its sign function must have exactly $n$ sign changes,
for otherwise the first $n$ moments already lie on $\partial M_n$
by the result just proved for $n-1$.  Such a sign function ends
with $+$; conversely every step function with exactly $n$ changes
and ending with $+$ is exposed by a degree-$n$ polynomial with
positive leading coefficient and therefore belongs to this upper
endpoint graph.  The lower endpoint graph similarly consists in
its interior of the functions with exactly $n$ changes ending
with $-$.

For $x\in\partial M_n$, the whole fiber is a singleton: a support
normal $c$ of $M_n$ at $x$ exposes a single point of $M_{n+1}$
with normal $(c,0)$, and every point of the fiber lies in that
support face.  The endpoint functions $u,\ell$ are continuous.
They are respectively concave and convex and hence continuous on
the interior of the convex body.  At a boundary point, compactness
of $M_{n+1}$ and the singleton-fiber assertion force every convergent
subsequence of endpoint values to have the unique boundary value;
boundedness then gives continuity there as well.  Thus each endpoint
graph is homeomorphic by projection to the closed ball $M_n$, and
their common boundary is the image of $|K_{n-1}|$.

Finally the image of either alternating $n$-simplex consists of its
open top simplex together with all of $K_{n-1}$: its two end faces
are the two alternating $(n-1)$-simplices, and all lower simplices
occur recursively.  We have therefore identified exactly that image
with one of the two endpoint graphs.  For $n=0$ the image is a
single point, which supplies the initial case.

\paragraph{Consequence for a removed lower simplicial subset.}
If $A$ is the image of an alternating $n$-simplex and $L\subset A$
is a simplicial subset omitting its top nondegenerate simplex, then
$|L|$ is a closed subset of the boundary of this closed ball.
After a homeomorphism with a convex closed ball, straight-line
contraction to an interior point stays outside $|L|$ at every
positive time.  Hence $|A|\setminus|L|$ is contractible, including
the case $L=\partial A$.  This proves the geometric ingredient of
B.4 without requiring $L$ itself to be a ball.

\paragraph{The relative simplex-category comparison.}
For any simplicial set $A$ and simplicial subset $L\subseteq A$,
let $\mathcal C=\Delta/A$ and let $\mathcal C_{\mathrm{out}}$ be
its full subcategory on the simplices which do not belong to $L$.
There is a specified zigzag of weak equivalences
\[
 B\mathcal C_{\mathrm{out}}
 \ \longleftarrow\ \hocolim_{\mathcal C_{\mathrm{out}}}U
 \ \longrightarrow\ |A|\setminus|L|,
 \qquad
 U(\sigma)=|\Delta^m|\setminus|\sigma^{-1}L|
 \quad(\sigma:\Delta[m]\to A).
\]
We prove this for arbitrary small simplicial sets, including
$A=L$ and $L=\emptyset$.

First extend $U$ to all of $\mathcal C$ by the same formula;
$U(\sigma)$ is empty exactly for $\sigma\in L$.  If
$\alpha:\sigma\to\tau$ is a simplex-category arrow, its affine
map takes $U(\sigma)$ to $U(\tau)$ because
$\tau\alpha=\sigma$.  More precisely the preimage of $|L|$ in
$|\Delta^m|$ is exactly $|\sigma^{-1}L|$.  Indeed a point has a
unique expression in the interior of a nondegenerate simplex, and
it belongs to $|L|$ exactly when that nondegenerate simplex belongs
to $L$.  On a standard simplex this tests the face given by the
positive barycentric coordinates.  Thus the removed set is a union
of closed faces, lying in the boundary whenever $\sigma\notin L$.
Straight-line contraction to the barycenter proves that every
nonempty $U(\sigma)$ is contractible.

The category $\mathcal C$ is Reedy: degree is simplex dimension,
direct arrows are injective operators, inverse arrows are surjective
operators, and the unique image factorization of an order-preserving
map supplies the Reedy factorization.  Its inverse matching category
at a nondegenerate simplex is empty.  At a degenerate simplex it
has terminal object the unique surjection onto its nondegenerate
core: every other degeneracy representation factors uniquely through
this one by the Eilenberg--Zilber decomposition.  Hence $\mathcal C$
has fibrant constants.  These are precisely Hirschhorn's
Example~15.1.19 and Proposition~15.10.4(1), with the proof just
checked; Theorem~15.10.8(2) says that colimit is then left Quillen
for the Reedy model structure \cite{BHirschhorn2003}.

The diagram $U$ is Reedy cofibrant in the usual model category of
compactly generated spaces.  Its latching map at $\sigma$ is
\[
 |\partial\Delta^m|\setminus|\sigma^{-1}L|
    \ \lhook\joinrel\longrightarrow\
 |\Delta^m|\setminus|\sigma^{-1}L|.
\]
To verify an actual cofibration, rather than merely a homotopy-type
assertion, put $Z=|\sigma^{-1}L|$.  If $Z$ is empty we have the
usual finite CW pair, and if $Z$ is the whole simplex both spaces
are empty.  Otherwise use barycentric coordinates and the continuous
piecewise linear function
\[
 \delta(x)=
 \max_{\substack{\emptyset\ne S\subseteq\{0,\ldots,m\}\\
                  \Delta^S\not\subseteq Z}}
       \min_{i\in S}x_i.
\]
It vanishes exactly on $Z$: a support face belongs to $Z$ if and
only if every nonempty subset of that support belongs to $Z$.
The dyadic shells
$2^{-k-1}\leq\delta\leq2^{-k}$, together with the uppermost
shell, are finite polyhedra.  Cutting by all simplex faces and
these level sets and taking compatible barycentric subdivisions
triangulates the complement locally finitely, because every point
of the complement has $\delta>0$.  In this triangulation its
boundary complement is a subcomplex.  Thus the displayed map is
a relative CW inclusion.  The finite union of proper faces is
exactly the latching object, with its subspace topology; this also
follows by gluing its finite closed face cover.  Hence all latching
maps are cofibrations, as asserted.

Hirschhorn's Theorem~19.9.1(1), applied with the checked hypotheses,
now gives a weak equivalence
\[
 \hocolim_{\mathcal C}U\longrightarrow\colim_{\mathcal C}U.
\]
The colimit on the right is precisely $|A|\setminus|L|$.
The usual realization quotient has $U(\sigma)$ as the preimage
of this open complement on each simplex, and restricting a quotient
map to the inverse image of an open subset is a quotient map.
This proves both the set and topology assertions, with no finiteness
restriction on $A$.

Finally $\Delta/L\subseteq\mathcal C$ is a sieve.  In the bar
formula for this homotopy colimit, a nonempty summand is indexed
by a chain $\sigma_0\to\cdots\to\sigma_q$ with
$U(\sigma_0)\ne\emptyset$.  Such a chain cannot enter the sieve,
so it lies entirely in $\mathcal C_{\mathrm{out}}$.  Hence the
homotopy colimit over $\mathcal C$ is exactly that over
$\mathcal C_{\mathrm{out}}$.  The objectwise contraction maps
$U\to *$ induce a weak equivalence to
$\hocolim_{\mathcal C_{\mathrm{out}}}*=B\mathcal C_{\mathrm{out}}$
by homotopy invariance, Hirschhorn's Theorem~19.4.2(1).
This supplies the announced zigzag.  For $A=L$ all three spaces
are empty, so that edge case is included.
